%============================================================
% CH2 · P3 · PROBLEM 3 — MAZE (CF 377A)
%============================================================
\subsection{Problem 2.3.3: Maze}
\problemlabel{2.3.3}{ssec:ch2-maze}

\problemcard{2.3.3}{Maze}{Codeforces 377A}{https://codeforces.com/problemset/problem/377/A}{\exmedium\ \cf{1600}}{Wall up the first $k$ cells of the DFS finishing order (they form whole subtrees), or keep the first $s-k$ cells of any traversal order (a connected prefix)}

\tikzset{
    g2p3z-e/.style={gcell, minimum size=6.5mm, font=\scriptsize\bfseries},
    g2p3z-w/.style={gwall, minimum size=6.5mm},
    g2p3z-x/.style={gcell, minimum size=6.5mm, fill=dupfill, draw=recursionred, font=\scriptsize\bfseries, text=recursionred},
    g2p3z-c/.style={gcell, minimum size=6.5mm, fill=currentfill, draw=accentorange, line width=1pt, font=\scriptsize\bfseries},
    g2p3z-s/.style={gcell, minimum size=6.5mm, fill=hitfill, font=\scriptsize\bfseries},
    g2p3z-d/.style={gcell, minimum size=6.5mm, fill=computedfill, font=\scriptsize\bfseries},
}
% \gtpz{row}{col}{style}{text}: one cell of a maze drawing (row 0 on top, 0.65 cm per cell)
\newcommand{\gtpz}[4]{\node[#3] at (#2*0.65, -#1*0.65) {#4};}

\subsubsection*{1. Problem statement}

A maze is an $n \times m$ grid of empty cells (\texttt{.}) and walls (\texttt{\#}); moves go between side-adjacent empty cells, and all empty cells form one connected area. Turn \textbf{exactly} $k$ empty cells into walls, marked \texttt{X}, so that the remaining empty cells are still connected. Print the new maze; a solution is guaranteed and any valid one is accepted.

\textbf{Constraints:} $1 \le n, m \le 500$, $0 \le k < s$ where $s$ is the number of empty cells.

\begin{table}[H]
\centering\small
\begin{tabular}{@{}p{5.2cm}p{4.3cm}p{4.3cm}@{}}
\toprule
\textbf{Input} & \textbf{Official output} & \textbf{Our output} \\
\midrule
\texttt{3 4 2 / \#..\# / ..\#. / \#...} & \texttt{\#.X\# / X.\#. / \#...} & \texttt{\#..\# / ..\#X / \#..X} \\
\texttt{5 4 5 / \#... / \#.\#. / .\#.. / ...\# / .\#.\#} & \texttt{\#XXX / \#X\#. / X\#.. / ...\# / .\#.\#} & \texttt{\#... / \#X\#. / X\#.. / X..\# / X\#X\#} \\
\bottomrule
\end{tabular}
\end{table}

\begin{diagrambox}
\begin{tikzpicture}
    \begin{scope}
        \foreach \r/\c/\s in {0/0/w, 0/1/e, 0/2/e, 0/3/w, 1/0/e, 1/1/e, 1/2/w, 1/3/e, 2/0/w, 2/1/e, 2/2/e, 2/3/e}
            \gtpz{\r}{\c}{g2p3z-\s}{}
        \node[font=\scriptsize\bfseries] at (0.98,0.6) {sample 1, $k = 2$};
    \end{scope}
    \begin{scope}[xshift=3.7cm]
        \foreach \r/\c/\s/\t in {0/0/w/, 0/1/e/, 0/2/x/X, 0/3/w/, 1/0/x/X, 1/1/e/, 1/2/w/, 1/3/e/, 2/0/w/, 2/1/e/, 2/2/e/, 2/3/e/}
            \gtpz{\r}{\c}{g2p3z-\s}{\t}
        \node[font=\scriptsize\bfseries] at (0.98,0.6) {official answer};
    \end{scope}
    \begin{scope}[xshift=7.4cm]
        \foreach \r/\c/\s/\t in {0/0/w/, 0/1/e/, 0/2/e/, 0/3/w/, 1/0/e/, 1/1/e/, 1/2/w/, 1/3/x/X, 2/0/w/, 2/1/e/, 2/2/e/, 2/3/x/X}
            \gtpz{\r}{\c}{g2p3z-\s}{\t}
        \node[font=\scriptsize\bfseries] at (0.98,0.6) {our answer};
    \end{scope}
    \node[note, text width=3.6cm, anchor=west, align=left] at (10.6,-0.65) {two different answers, both valid: the 6 remaining empty cells are connected in each};
\end{tikzpicture}
\figcap{Sample 1. Dark cells are walls, red cells the $k$ new walls \texttt{X}. The official answer removes two dead ends at the top; ours removes the tail $(1,3), (2,3)$ on the right.}
\end{diagrambox}

\subsubsection*{2. The graph model}

\begin{modelbox}[Graph Model --- Maze]
\renewcommand{\arraystretch}{1.35}
\begin{tabularx}{\linewidth}{@{}>{\bfseries}l X@{}}
Vertices & the $s$ empty cells \\
Edges & between side-adjacent empty cells, implicit (computed from the grid) \\
Weights & none \\
Question & delete exactly $k$ vertices of a connected graph so that the rest stays connected --- equivalently, choose $s - k$ vertices inducing a connected subgraph \\
Size & $s \le 2.5\cdot10^5$, $|E| < 5\cdot10^5$: one traversal; a DFS can be $1.25\cdot10^5$ cells deep \\
\end{tabularx}
\end{modelbox}

\begin{diagrambox}
\begin{tikzpicture}[x=1.1cm, y=1.1cm]
    \tikzset{gv/.append style={minimum size=8mm, font=\tiny\bfseries}}
    \foreach \r/\c in {0/1, 0/2, 1/0, 1/1, 1/3, 2/1, 2/2, 2/3}
        \node[gv] (g\r\c) at (\c,-\r) {\r,\c};
    \foreach \a/\b in {01/02, 01/11, 10/11, 11/21, 21/22, 22/23, 13/23} \draw[ge] (g\a) -- (g\b);
    \node[note, anchor=west, text width=5.6cm, align=left] at (4.0,-1.0) {the 8 empty cells of sample 1 as a graph (labels: row, column). Deleting a vertex keeps it connected iff the vertex is not a cut vertex; leaves of any spanning tree are never cut vertices.};
\end{tikzpicture}
\figcap{The model: cells are vertices, shared sides are edges. Here the graph happens to be a tree, so the safe cells to remove first are exactly its leaves $(0,2)$, $(1,0)$, $(1,3)$.}
\end{diagrambox}

\subsubsection*{3. Key observations}

\begin{itemize}[leftmargin=*]
    \item \textbf{Greedy removal one by one works} --- a connected graph always has a vertex whose removal keeps it connected (a leaf of a spanning tree) --- but finding such a vertex anew $k$ times is slow.
    \item \textbf{All $k$ at once with the finishing order.} Every prefix of the DFS post-order is a union of complete DFS subtrees; the unfinished cells always contain the root and the tree path from each of them to the root (invariant of the pattern). So the first $k$ cells to finish can be walled up together.
    \item \textbf{The mirror view: keep a connected prefix.} In a BFS (or DFS pre-) order every cell except the start is discovered from an earlier cell, so the first $s - k$ cells are connected. Walling up the \emph{last} $k$ cells of a BFS order is equally correct.
    \item \textbf{Exactly $k$, not at most.} Both rules produce exactly $k$ walls, since $k < s$.
    \item \textbf{Depth.} A snake-shaped $500 \times 500$ maze has $125{,}250$ empty cells in one path: the DFS must be iterative (or use BFS).
\end{itemize}

\subsubsection*{4. Approaches}

\begin{bruteforcebox}
\textbf{Idea: remove one safe cell at a time.} $k$ times, scan the empty cells in row-major order and wall up the first one whose removal leaves the other empty cells connected (checked by a BFS).

\begin{lstlisting}[style=cpp]
#include <bits/stdc++.h>
using namespace std;

// Brute force: k times, try the empty cells in row-major order and wall up the first one
// whose removal leaves the other empty cells connected (checked by a BFS).
// A connected graph always has such a cell (a leaf of any spanning tree), so this never gets stuck.
int n, m;
vector<string> g;

bool connectedWithout(int skip) {
    int start = -1, total = 0;
    for (int i = 0; i < n * m; i++)
        if (i != skip && g[i / m][i % m] == '.') { total++; if (start < 0) start = i; }
    if (total == 0) return true;
    vector<char> seen(n * m, 0);
    queue<int> q;
    q.push(start); seen[start] = 1;
    int reached = 0;
    const int dr[4] = {-1, 1, 0, 0}, dc[4] = {0, 0, -1, 1};
    while (!q.empty()) {
        int c = q.front(); q.pop(); reached++;
        for (int d = 0; d < 4; d++) {
            int nr = c / m + dr[d], nc = c % m + dc[d];
            if (nr < 0 || nr >= n || nc < 0 || nc >= m) continue;
            int x = nr * m + nc;
            if (x == skip || g[nr][nc] != '.' || seen[x]) continue;
            seen[x] = 1; q.push(x);
        }
    }
    return reached == total;
}

int main() {
    int k;
    cin >> n >> m >> k;
    g.resize(n);
    for (auto& row : g) cin >> row;
    for (int t = 0; t < k; t++)
        for (int i = 0; i < n * m; i++)
            if (g[i / m][i % m] == '.' && connectedWithout(i)) { g[i / m][i % m] = 'X'; break; }
    for (auto& row : g) cout << row << "\n";
}
\end{lstlisting}
On sample 1 it walls up $(0,2)$ (removing $(0,1)$ first would isolate $(0,2)$) and then $(0,1)$, printing \texttt{\#XX\# / ..\#. / \#...} --- a third valid answer.

\begin{complexitybox}
\textbf{Time:} up to $k \cdot s$ connectivity checks of \bigO{nm} each; for an open $500 \times 500$ grid that is far beyond $10^{15}$. \textbf{Space:} \bigO{nm}. \textbf{Verdict:} \tle{}; the reference for random tests (together with a checker, since answers are not unique).
\end{complexitybox}
\end{bruteforcebox}

\begin{algorithmbox}[Approach 2: wall up the first $k$ cells to finish (iterative DFS)]
\textbf{Idea.} DFS from any empty cell. Each stack entry remembers the next of its four directions to try (\texttt{dir[c]}); when all four are tried, the cell \emph{finishes}: if fewer than $k$ cells have been walled so far, mark it \texttt{X}.

\begin{lstlisting}[style=cpp]
#include <bits/stdc++.h>
using namespace std;

// Maze (CF 377A): wall up the FIRST k cells in DFS finishing order.
// A finished set is always a union of whole DFS subtrees, so the unfinished rest stays connected.
// Iterative: a snake-shaped maze of 500 x 500 makes the DFS 125,000+ cells deep.
int main() {
    ios::sync_with_stdio(false);
    cin.tie(nullptr);
    int n, m, k;
    cin >> n >> m >> k;
    vector<string> g(n);
    for (auto& row : g) cin >> row;
    int start = -1;
    for (int i = 0; i < n * m && start < 0; i++)
        if (g[i / m][i % m] == '.') start = i;
    const int dr[4] = {-1, 1, 0, 0}, dc[4] = {0, 0, -1, 1};
    vector<char> seen(n * m, 0), dir(n * m, 0);      // dir[c] = next direction to try from c
    vector<int> st(1, start);
    seen[start] = 1;
    int walled = 0;
    while (!st.empty()) {
        int c = st.back(), r = c / m, col = c % m;
        if (dir[c] == 4) {                           // c finishes
            st.pop_back();
            if (walled < k) { g[r][col] = 'X'; walled++; }
            continue;
        }
        int d = dir[c]++;
        int nr = r + dr[d], nc = col + dc[d];
        if (nr < 0 || nr >= n || nc < 0 || nc >= m || g[nr][nc] != '.' || seen[nr * m + nc]) continue;
        seen[nr * m + nc] = 1;
        st.push_back(nr * m + nc);
    }
    for (auto& row : g) cout << row << "\n";
}
\end{lstlisting}

\begin{complexitybox}
\textbf{Time:} \bigO{nm}: each cell is pushed once and tries four directions. \textbf{Space:} \bigO{nm}. \textbf{Verdict:} \ac{}-ready: about $0.2$\,s here on the $500 \times 500$ snake ($125{,}250$ cells, DFS depth $125{,}250$) and on the open $500 \times 500$ grid with $k = s - 1$; the checker accepts both outputs.
\end{complexitybox}
\end{algorithmbox}

\begin{invariantbox}[Why the remaining cells stay connected]
Let $F$ be the set of finished cells at any moment. A cell finishes only after every cell it discovered has finished, so $F$ is closed under taking DFS-tree descendants. Take any unfinished cell $v$ other than the start: its DFS parent finishes after $v$, so it is unfinished too. Repeating, the whole tree path from $v$ to the start avoids $F$. Hence every unfinished cell is joined to the start inside the unfinished set.

The program walls up exactly the cells that finish while \texttt{walled < k}, i.e.\ the first $k$ cells of the finishing order --- the set $F$ at the moment the $k$-th cell finished. The remaining $s-k \ge 1$ cells are unfinished at that moment and therefore connected. Writing \texttt{X} during the search does not disturb it: a cell is marked only after it has finished, and a finished cell is already \texttt{seen}, so it is never pushed again. \hfill$\blacksquare$
\end{invariantbox}

\begin{alternativebox}[Approach 3: keep a connected prefix --- wall up the last $k$ cells of a BFS order]
The other side of the same coin: the first $j$ cells of a BFS order are connected for every $j$ (each was discovered from an earlier one). Keep the first $s - k$, wall up the rest. BFS needs no direction counters and no deep stack; the queue is the order.

\begin{lstlisting}[style=cpp]
#include <bits/stdc++.h>
using namespace std;

// Alternative: keep a connected PREFIX. In BFS order every cell (except the start) is
// discovered from an earlier cell, so the first s - k cells are connected: wall up the last k.
int main() {
    ios::sync_with_stdio(false);
    cin.tie(nullptr);
    int n, m, k;
    cin >> n >> m >> k;
    vector<string> g(n);
    for (auto& row : g) cin >> row;
    int start = -1;
    for (int i = 0; i < n * m && start < 0; i++)
        if (g[i / m][i % m] == '.') start = i;
    const int dr[4] = {-1, 1, 0, 0}, dc[4] = {0, 0, -1, 1};
    vector<char> seen(n * m, 0);
    vector<int> order;                               // BFS order = the queue itself
    order.push_back(start);
    seen[start] = 1;
    for (size_t h = 0; h < order.size(); h++) {
        int r = order[h] / m, c = order[h] % m;
        for (int d = 0; d < 4; d++) {
            int nr = r + dr[d], nc = c + dc[d];
            if (nr < 0 || nr >= n || nc < 0 || nc >= m || g[nr][nc] != '.' || seen[nr * m + nc]) continue;
            seen[nr * m + nc] = 1;
            order.push_back(nr * m + nc);
        }
    }
    int s = order.size();
    for (int i = s - k; i < s; i++) g[order[i] / m][order[i] % m] = 'X';
    for (auto& row : g) cout << row << "\n";
}
\end{lstlisting}

\begin{diagrambox}
\begin{tikzpicture}
    \begin{scope}
        \foreach \r/\c/\s/\t in {0/0/w/, 0/1/d/0, 0/2/d/2, 0/3/w/, 1/0/d/4, 1/1/d/1, 1/2/w/, 1/3/x/7, 2/0/w/, 2/1/d/3, 2/2/d/5, 2/3/x/6}
            \gtpz{\r}{\c}{g2p3z-\s}{\t}
        \node[font=\scriptsize\bfseries] at (0.98,0.6) {BFS order, keep $0..5$};
    \end{scope}
    \begin{scope}[xshift=4.2cm]
        \foreach \r/\c/\s/\t in {0/0/w/, 0/1/d/8, 0/2/d/7, 0/3/w/, 1/0/d/5, 1/1/d/6, 1/2/w/, 1/3/x/1, 2/0/w/, 2/1/d/4, 2/2/d/3, 2/3/x/2}
            \gtpz{\r}{\c}{g2p3z-\s}{\t}
        \node[font=\scriptsize\bfseries] at (0.98,0.6) {DFS finish \#, remove 1, 2};
    \end{scope}
    \node[note, text width=5.6cm, anchor=west, align=left] at (7.6,-0.65) {On sample 1 both rules wall up the same two cells; on sample 2 they differ (\texttt{\#... / \#X\#. / X\#.. / X..\# / X\#X\#} by DFS, \texttt{\#... / \#.\#. / X\#.. / XX.\# / X\#X\#} by BFS), and both pass the checker.};
\end{tikzpicture}
\figcap{The two orders on sample 1. Left: BFS positions; the last $k = 2$ become walls. Right: DFS finishing numbers; the first $k = 2$ become walls.}
\end{diagrambox}

\begin{complexitybox}
\textbf{Time:} \bigO{nm}. \textbf{Space:} \bigO{nm}. \textbf{Verdict:} \ac{}-ready, also about $0.2$\,s on the two $500 \times 500$ tests. The DFS \emph{pre}-order would work just as well as BFS for ``keep a prefix''; the post-order is the natural order for ``remove a prefix''.
\end{complexitybox}
\end{alternativebox}

\begin{tricksbox}[Approach 4: Implementation Refinements]
\begin{itemize}[leftmargin=*]
    \item \textbf{Recursive DFS is a trap here.} The snake maze sends a recursive DFS $125{,}250$ frames deep --- about twice what a 2\,MB stack survives (Lesson~\ref{ssec:ch2-call-stack}). Either the iterative form above, or BFS.
    \item \textbf{One index per cell.} \texttt{c = r * m + col} keeps the stack and the arrays one-dimensional; \texttt{dir} fits in a \texttt{char}.
    \item \textbf{Mark \texttt{X} in place.} After a cell finishes it is never examined again, so writing \texttt{X} immediately is safe; the BFS version writes after the search.
    \item \textbf{Start anywhere.} Any empty cell can be the root; the first in row-major order is simplest.
\end{itemize}
\end{tricksbox}

\subsubsection*{5. Trace}

\begin{tracebox}[{Sample 1, directions tried up, down, left, right; start $(0,1)$, $k = 2$}]
\begin{center}
\begin{tikzpicture}
    \begin{scope}
        \foreach \r/\c/\s/\t in {0/0/w/, 0/1/s/0, 0/2/e/, 0/3/w/, 1/0/e/, 1/1/s/1, 1/2/w/, 1/3/c/5, 2/0/w/, 2/1/s/2, 2/2/s/3, 2/3/s/4}
            \gtpz{\r}{\c}{g2p3z-\s}{\t}
        \node[font=\scriptsize\bfseries] at (0.98,0.6) {enter $(1,3)$, depth 5};
        \draw[primaryblue, line width=1.4pt, -{Stealth[length=2mm]}] (0.65,-0.2) -- (0.65,-0.45);
        \draw[primaryblue, line width=1.4pt, -{Stealth[length=2mm]}] (0.65,-0.85) -- (0.65,-1.1);
        \draw[primaryblue, line width=1.4pt, -{Stealth[length=2mm]}] (0.85,-1.3) -- (1.1,-1.3);
        \draw[primaryblue, line width=1.4pt, -{Stealth[length=2mm]}] (1.5,-1.3) -- (1.75,-1.3);
        \draw[primaryblue, line width=1.4pt, -{Stealth[length=2mm]}] (1.95,-1.1) -- (1.95,-0.85);
    \end{scope}
    \begin{scope}[xshift=3.9cm]
        \foreach \r/\c/\s/\t in {0/0/w/, 0/1/s/, 0/2/e/, 0/3/w/, 1/0/e/, 1/1/s/, 1/2/w/, 1/3/x/X, 2/0/w/, 2/1/s/, 2/2/c/, 2/3/x/X}
            \gtpz{\r}{\c}{g2p3z-\s}{\t}
        \node[font=\scriptsize\bfseries] at (0.98,0.6) {2 finished: $k$ reached};
    \end{scope}
    \begin{scope}[xshift=7.8cm]
        \foreach \r/\c/\s/\t in {0/0/w/, 0/1/d/, 0/2/d/, 0/3/w/, 1/0/d/, 1/1/d/, 1/2/w/, 1/3/x/X, 2/0/w/, 2/1/d/, 2/2/d/, 2/3/x/X}
            \gtpz{\r}{\c}{g2p3z-\s}{\t}
        \node[font=\scriptsize\bfseries] at (0.98,0.6) {end: 6 cells kept};
    \end{scope}
    \node[note, text width=3.4cm, anchor=west, align=left] at (11.0,-0.65) {orange $=$ top of stack, light orange $=$ on the stack, blue $=$ finished and kept, red $=$ walled};
\end{tikzpicture}
\end{center}
\begin{center}\small
\begin{tabular}{@{}clp{5.6cm}l@{}}
\toprule
\textbf{\#} & \textbf{event} & \textbf{cell (entered from)} & \textbf{action} \\
\midrule
1 & enter 0 & $(0,1)$ & start \\
2 & enter 1--5 & $(1,1), (2,1), (2,2), (2,3), (1,3)$, each from the previous & down, down, right, right, up \\
3 & finish 1 & $(1,3)$ & \texttt{X} (walled 1) \\
4 & finish 2 & $(2,3)$ & \texttt{X} (walled 2 $= k$) \\
5 & finish 3, 4 & $(2,2), (2,1)$ & kept \\
6 & enter 6, finish 5 & $(1,0)$ (from $(1,1)$) & kept \\
7 & finish 6 & $(1,1)$ & kept \\
8 & enter 7, finish 7, 8 & $(0,2)$ (from $(0,1)$), then $(0,1)$ & kept \\
\bottomrule
\end{tabular}
\end{center}
Finishing order: $(1,3), (2,3), (2,2), (2,1), (1,0), (1,1), (0,2), (0,1)$. The first two are a whole subtree (the tail hanging below $(2,2)$).
\end{tracebox}

\subsubsection*{6. Test suite}

\begin{testbox}
\centering\small
\begin{tabular}{@{}p{4.8cm}p{3.6cm}p{5.2cm}@{}}
\toprule
\textbf{Input} & \textbf{Our output} & \textbf{Why} \\
\midrule
\texttt{1 1 0 / .} & \texttt{.} & one cell, nothing to do \\
\texttt{1 5 3 / .....} & \texttt{..XXX} & a path: the far end finishes first \\
\texttt{2 2 3 / .. / ..} & \texttt{.X / XX} & a 4-cycle, $k = s - 1$ \\
\texttt{3 3 4 / ... / .\#. / ...} & \texttt{.XX / .\#X / ..X} & a ring: removing a contiguous arc \\
$500 \times 500$ snake, $k = 62625$ & valid (checker) & DFS depth $125{,}250$ \\
$500 \times 500$ open, $k = s - 1$ & valid (checker) & only the start survives \\
\bottomrule
\end{tabular}
\end{testbox}
All outputs come from running Approach~2 (the BFS version prints \texttt{..X / .\#X / .XX} on the ring). Because answers are not unique, a checker verifies every output: exactly $k$ cells changed, each from \texttt{.} to \texttt{X}, nothing else changed, and the remaining \texttt{.} cells connected. Approaches 1--3 pass it on both official samples and on $1000$ random mazes up to $6 \times 6$ (random walls, then only the largest component kept, random $k < s$).

\subsubsection*{7. Common mistakes}

\begin{warningbox}[Mistake 1 --- walling the first $k$ cells \emph{entered}]
Pre-order is the order for keeping, not for removing. On sample 1 the first two cells entered are $(0,1)$ and $(1,1)$; walling them up cuts $(0,2)$ off from everything else.
\end{warningbox}

\begin{warningbox}[Mistake 2 --- any $k$ cells with degree 1]
``Remove dead ends'' is right one cell at a time, but $k$ may exceed the number of dead ends (on an open grid there are none), and removing one dead end creates new ones. The finishing order handles all of that in one pass.
\end{warningbox}

\begin{warningbox}[Mistake 3 --- recursion]
The snake maze is a legal test; recursive DFS crashes at depth $125{,}250$.
\end{warningbox}

\begin{warningbox}[Mistake 4 --- changing walls or printing \texttt{\#}]
Only empty cells may change, and they must be printed as \texttt{X}, not \texttt{\#}. A checker rejects both.
\end{warningbox}

\subsubsection*{8. Takeaways}

\begin{itemize}[leftmargin=*]
    \item Every prefix of a DFS finishing order is a union of complete subtrees: removing it keeps the rest connected.
    \item Every prefix of a BFS (or DFS entry) order is connected: keeping it is the mirror statement.
    \item ``Exactly $k$, any valid answer'' asks for a construction, and a traversal order is often the construction.
    \item Grids up to $500 \times 500$ need an iterative DFS or a BFS.
\end{itemize}

\begin{variantsbox}
\variant{Keep exactly $t$ cells}{\varnone}
  {keep a connected region of exactly $t$ cells that contains a given cell $c$.}
  {start the traversal at $c$.}
  {BFS from $c$, keep the first $t$ cells (or remove the first $s - t$ finishers of a DFS rooted at $c$, which never removes $c$).}
\variant{Remove vertices of a general graph}{\varnone}
  {a connected graph (not a grid); delete $k$ vertices keeping the rest connected.}
  {same model, explicit adjacency lists.}
  {the first $k$ vertices of the finishing order of a DFS (template \texttt{Timestamps::post} or \texttt{finishOrder}).}
\variant{Dismantle a tree leaf by leaf}{\varnone}
  {print an order in which the vertices of a tree can be deleted so that the remaining tree stays connected after every step.}
  {the tree itself.}
  {the post-order of any root is such an order (each prefix is a union of subtrees); the reverse of a BFS order works too.}
\variant{Topological order of a DAG}{\varnone}
  {order the vertices of a DAG so that all edges point forwards.}
  {the finishing order again, now in a directed graph.}
  {reverse the finishing order (pattern template \texttt{finishOrder}; Chapter~5).}
\end{variantsbox}
